\chapter[Band 31: Rechtecksbänder]{Band 31 auf einen Blick}
\noindent\textbf{Rechtecksbänder}\par\medskip
Vorausgesetzt sind Halbgruppen, idempotente Operationen und
Potenzhalbgruppen. \(\mathcal P_+(A)\) bezeichnet die nichtleeren Teilmengen
mit Mengenprodukt; \(\cong\) stets Halbgruppenisomorphie.

\begin{BandUebersicht}
\textbf{Rechtecksidentität}
\UebFormel{\begin{aligned}
&\RectangularBandAlg R\star\ \leftrightarrow\\
&\quad\SemigroupAlg R\star\land\\
&\quad\forall x,y\in R\ ((x\star y)\star x=x).
\end{aligned}}
\UebQuellen{\FormulaRefAuto{RectangularBandDef}[def]}
&
\textbf{Idempotenz und Dreigliedrigkeit}
\UebFormel{\begin{aligned}
&\RectangularBandAlg R\star\ \vdash\
 \IdempotentSemigroupAlg R\star,\\
&\RectangularBandAlg R\star,\ x,y,z\in R\ \vdash\\
&\qquad(x\star y)\star z=x\star z.
\end{aligned}}
\UebQuellen{\FormulaRefAuto{RectangularBandIsIdempotentSemigroup}[thm];
\FormulaRefAuto{RectangularBandThreeTermLaw}[thm]} \\
\addlinespace[.8ex]

\textbf{Produktdarstellung}
Für \(\SemigroupAlg R\star\) und \(e\in R\) sind die Koordinaten\UebFormel{\begin{aligned}
&\mathcal K_e=(Re)\times(eR),\\
&(p,q)\square_e(r,s)=(p,s),\\
&\chi_e(x)=(x\star e,e\star x).
\end{aligned}}
\UebQuellen{\FormulaRefAuto{SemigroupCoordinateCarrierDef}[def];
\FormulaRefAuto{SemigroupCoordinateProductDef}[def];
\FormulaRefAuto{SemigroupCoordinateMapDef}[def]}
&
\textbf{Charakterisierung nichtleerer Rechtecksbänder}
\UebFormel{\begin{aligned}
&R\ne\varnothing\ \vdash\\
&\bigl(\RectangularBandAlg R\star\ \leftrightarrow\\
&\qquad\forall e\in R\
 \SemigroupCoordinateIso R\star e\bigr).
\end{aligned}}
\UebQuellen{\FormulaRefAuto{RectangularBandProductRepresentation}[thm]} \\
\addlinespace[.8ex]

\textbf{Rekonstruktionsproblem für endliche Bänder}
Betrachtet werden endliche idempotente Halbgruppen
\((A,\star)\), \((B,\diamond)\). Eine Potenzhalbgruppenisomorphie
ist eine Bijektion \(\Psi:\mathcal P_+(A)\to\mathcal P_+(B)\) mit\UebFormel{\begin{aligned}
&\Psi(X\star_{\mathcal P}Y)
   =\Psi(X)\diamond_{\mathcal P}\Psi(Y).
\end{aligned}}
\UebQuellen{\FormulaRefAuto{IdempotentSemigroupDef}[def];
\FormulaRefAuto{SemigroupIsoDef}[def]}
&
\textbf{Bestimmtheit in der bewiesenen endlichen Klasse}
\UebFormel{\begin{aligned}
&\IdempotentSemigroupAlg A\star,\
 \IdempotentSemigroupAlg B\diamond,\\
&\Finite A,\ \Finite B,\
 \mathcal P_+(A)\cong\mathcal P_+(B)\ \vdash\\
&\qquad A\cong B.
\end{aligned}}Der Band führt dies über Bandzerlegung, Komponententransport
und endliches Anheben auf die Grundelemente. Die Endlichkeits- und
Idempotenzvoraussetzungen bleiben bestehen.
\UebQuellen{\FormulaRefAuto{FiniteIdempotentSemigroupsGloballyDetermined}[thm]} \\
\addlinespace[.8ex]
\end{BandUebersicht}

